\section{Library Learning modulo Equational Theory}\label{sec:egraphs}

\begin{figure}
  \begin{minipage}{.5\textwidth}
    \textbf{Syntax}
    $$
    \begin{array}{rll}
      \text{e-class ids} & a,b &\in \mathcal{I} \\
      \text{e-nodes}     & n \gramdef s(a_1, \ldots, a_k)  &\in N\\
      \text{e-classes}   & c \gramdef \{n_1, \ldots, n_m\} &\in C     
    \end{array}
    $$      
  \end{minipage}%
  \begin{minipage}{.5\textwidth}
    \textbf{Denotation} $\denot{\cdot}\colon N\to 2^{\termsof{\Sigma}}$, $\denot{\cdot}\colon C\to 2^{\termsof{\Sigma}}$
    $$
    \begin{array}{rl}
      \\
      \denot{s(a_1, \ldots, a_k)} &= \{s(t_1, \ldots, t_k) \mid t_i \in \denot{M(a_i)}\}\\
      \denot{\{n_1, \ldots, n_m\}} &= \bigcup_{i \in [1,m]} \denot{n_i}
    \end{array}
    $$      
  \end{minipage}
  \caption{Syntax, metavariables, and denotation for the components of an e-graph.
           Here $s \in \Sigma$ and $k = \arity(s)$.}\label{fig:syntax}
\end{figure}

% \subsection{Preliminaries}

\mypara{E-Graphs}
%
Let $\mathcal{I}$ be a denumerable set of \emph{e-class ids}.
%
An \emph{e-graph} $\mathcal{G}$ is a triple $\langle C, M, r\rangle$,
where $C$ is a set of \emph{e-classes},
$M\colon \mathcal{I} \to C$ is an \emph{e-class map}.
and $r\in \mathcal{I}$ is the root class id.
% \footnote{
%   Other formalizations of e-graphs do not feature a distinguished root class.
%   We include it, because denoting the e-graph to the set of
%   terms represented by its root e-class (as opposed to other e-classes)
%   is useful for our purposes.
% }
%
An \emph{e-class} $c \in C$ is a set of \emph{e-nodes} $n \in N$,
and an e-node is a constructor applied to e-class ids.
%
The syntax of e-classes and e-nodes is summarized in \autoref{fig:syntax} (left).
% $$
% c ::= \{n_1, \ldots, n_m\}\quad\quad\quad
% n ::= s(a_1, \ldots, a_k) \quad\text{where}\ s\in \Sigma, k = \arity(s), a_i \in \mathcal{I}\\
% $$
An e-graph has to satisfy the \emph{congruence invariant},
which states that the e-graph has no two identical e-nodes
(or alternatively, that all e-classes are disjoint).%
\footnote{In a real e-graph implementation, the definitions of e-graphs and the congruence invariant
are more involved, 
because efficient merging of e-classes requires introducing a non-trivial equivalence relation over e-class ids;
these details are irrelevant for our purposes.
Also, other formalizations of e-graphs do not feature a distinguished root e-class.} 

The \emph{denotation} of an \emph{e-graph}---the set of terms it represents---%
is the denotation of its root e-class $\denot{M(r)}$,
where the denotation of e-classes and e-nodes is defined mutually-recursively
in \autoref{fig:syntax} (right).
%
Note that the denotation can be infinite if the e-graph has cycles.
% $$
% \denot{\many{n}} = \bigcup_n \denot{n}\quad\quad\quad
% \denot{s(a_1, \ldots, a_k)} = \{s(t_1, \ldots, t_k) \mid t_i \in \denot{M(a_i)}\}
% $$
An e-graph induces an equivalence relation $\eqg$, 
where $t_1 \eqg t_2$ iff there exists an e-class $c \in C$ such that $t_1 \in \denot{c} \wedge t_2\in \denot{c}$.

E-graphs provide means to \emph{extract} the cheapest term from an e-class according to some cost function:
$\extract_{\cost}(a) = \argmin_{t\in \denot{M(a)}} \cost(t)$.
%
If $\cost$ is \emph{local},
 meaning that the cost of a term can be computed 
 from the costs of its immediate children,
% \todo{david: we don't necessarily mean monotonic---we're referring to a term
% costing at leasting as much as the SUM of its children}
extraction can be done efficiently by a greedy algorithm,
which recursively extracts the best term from each e-class.

\mypara{E-Matching}
%
E-matching is a generalization of pattern matching to e-graphs,
where matching an e-class $c$ against a pattern $p$ yields a set of \emph{e-class substitutions} 
$\theta\colon \mathcal{X}\to \mathcal{I}$
such that $\sapp{\theta}{p}$ is a ``subgraph'' of $c$.
%
To formalize this notion, we introduce \emph{partial terms} $\pi \in \patsof{\Sigma}{\mathcal{I}}$,
which are terms whose leaves can be e-class ids
(or, alternatively, patterns with e-class ids for variables).
%
The containment relation $\contains{\pi}{a}$ for some e-class id $a$ is defined as follows:
$$
\contains{a}{a} \quad\quad\quad
\contains{s(\pi_1, \ldots, \pi_k)}{a}\ \text{iff}\ s(a_1, \ldots, a_k) \in M(a) \wedge \contains{\pi_i}{a_i}
$$
% if either $\pi = a$ or $\pi$'s constructor occurs as an e-node $n$ in $M(a)$ 
% and all $\pi$'s children are recursively contained in $n$'s children.
%
With this definition, a pattern $p$ matches an e-class $a$, $a \sub p$, if there exists an e-class substitution $\theta$,
such that $\contains{\sapp{\theta}{p}}{a}$.
%
We denote the set of such substitutions as $\matches{a}{p}$.

\mypara{Rewriting and Equality Saturation}
%
Equality saturation (EqSat)~\cite{tate2009equality,willsey2021egg} 
takes as input a term $t$ and a set of equations that induce an equivalence relation $\semeq$,
and produces an e-graph $\mathcal{G}$ such that $\denot{\mathcal{G}} = \{t' \mid t \semeq t'\}$
and $t \semeq t'$ iff $t \eqg t'$.
%
The core idea of EqSat is to convert equations into rewrite rules
and apply them to the e-graph in a non-destructive way:
so that the original term and the rewritten terms are both represented in the same e-class.
%
Applying a rewrite rule $p_1 \Rightarrow p_2$ to an e-class $a$ works as follows:
for each $\theta \in \matches{a}{p_1}$, 
we obtain the rewritten partial term $\pi' = \sapp{\theta}{p_2}$
and then add this partial term to the same e-class $a$,
restoring the congruence invariant
(\ie merging e-classes that now have identical e-nodes).

\subsection{Top-Level Algorithm}

\begin{figure}
\begin{lstlisting}[
    commentstyle=\small\ttfamily\textcolor{gray},
    basicstyle=\small\ttfamily,
    xleftmargin=2.5em,
    numbers=left]
# Original term $t$, set of equational rewrite rules $\ruleset_\semeq$, maximum library size $N$
def LLMT($t$, $R_\semeq$, $N$):
    # EqSat phase:
    $\mathcal{G}$ = egraph($t$)                            # initialize with a single term $t$
    $\mathcal{G}$ = eqSat($\mathcal{G}, \ruleset_\semeq$)  # $\mathcal{G}$ represents all terms that are $\semeq t$

    # candidate generation phase:
    $\patset$ = AU($\mathcal{G}$)                          # generate candidate patterns by anti-unification
    $\ruleset_\kappa$ = $\{\kappa(p) \mid p \in \patset\}$ # construct a compression rule from every pattern
    $\mathcal{G}'$ = eqSat($\mathcal{G}, \ruleset_\kappa$) # $\mathcal{G}'$ represents all ways to compress $\mathcal{G}$

    # candidate selection phase:
    $\ruleset'$ = select_library($\mathcal{G}', N$)        # select the best $N$ rules from $\ruleset_\kappa$ using beam search
    $\mathcal{G}''$  = eqSat($\mathcal{G}, R'$)            # $\mathcal{G}''$ represents all ways to compress $\mathcal{G}$ using the optimal library          
    return extract($\mathcal{G}''$)                        # extract the smallest compressed term from $\mathcal{G}''$
\end{lstlisting}
  \caption{The top-level LLMT algorithm.}
\label{fig:algo-llmt}
\end{figure}

%\begin{figure}
%\begin{lstlisting}[
%    basicstyle=\normalsize\ttfamily,
%    xleftmargin=2.5em,
%    numbers=left]
%def anti-unify:
%    input: E: an e-graph containing the set of programs to anti-unify
%    output: R: a set of rewrites which introduces function definitions and applications
%               at the point of use
%
%    # TODO: i forgot how this works lmao
%\end{lstlisting}
%\caption{The antiunification algorithm of \babble.}
%\label{fig:algo-au}
%\end{figure}


We can formalize the problem of \emph{library learning modulo equational theory} (LLMT) as follows:
given a term $t$ and a set of equations that induce an equivalence relation $\semeq$,
the goal is to find a compressed term $\hat{t} \in \absof{\Sigma}{\mathcal{X}}$,
such that $\hat{t} \evals t' \semeq t$ (for some $t'$),
and $\hat{t}$ has a minimal size.

% This problem can be equivalently restated as library learning \emph{over an e-graph}:
% given an e-graph $\mathcal{G} = (C, M, r)$,
% the goal is to find a minimal-size compressed term $\hat{t} \in \absof{\Sigma}{\mathcal{X}}$,
% such that $\hat{t} \evals t$ for some $t \in \denot{\mathcal{G}}$.
% %
% If we can solve this problem, we can also solve the original LLMT problem,
% building the e-graph by applying EqSat to the input term and equations.

Our top-level algorithm \T{LLMT} is depicted in \autoref{fig:algo-llmt}.
%
This algorithm takes as input the original corpus $t$ 
and the equational theory, represented as a set of require rules $\ruleset_\semeq$
(another input to the algorithm is the maximum size of the library;
this parameter is introduced for the sake of efficiency, as we explain in \autoref{sec:beam}).
%
As the first step (lines 4--5),
\T{LLMT} applies EqSat to obtain an e-graph $\mathcal{G}$ 
that represents all terms $t'$ such that $t' \semeq t$.
%
This reduces the LLMT problem to library learning over an e-graph:
\ie the goal is to find a minimal-size compressed term $\hat{t}$,
such that $\hat{t} \evals t'$ for some $t' \in \denot{\mathcal{G}}$.
%
Similarly to \autoref{sec:au},
we take a pattern-based approach to this problem,
that is, we select a set $\patset$ of \emph{candidate patterns}
and then perform compression rewrites using these patterns.

The rest of the algorithm is split into two phases:
\emph{candidate generation} and \emph{candidate selection}.
%
Candidate generation (lines 8--10)
first computes the set of candidate patterns $\patset$ using the anti-unification mechanism extended to e-graphs.
%
Then it creates a compression rule ($\kappa$-rule, see \autoref{sec:au}) from each candidate pattern,
and once again applies EqSat to obtain a new e-graph $\mathcal{G}'$.
%
This new e-graph represents all possible ways to compress the terms from $\mathcal{G}$ 
using patterns in $\patset$.
%
Finally, the candidate selection phase (lines 13--15)
selects the optimal subset of compression rules $\ruleset'$,
constructs an e-graph $\mathcal{G}''$ that represents all possible compressions using only the selected compression rules,
and finally extracts the smallest compressed term from this e-graph.

The rest of this section focuses on the candidate generation via e-graph anti-unification
(line 8).
%
The candidate selection functions \T{select\_library} and \T{extract} are discussed in \autoref{sec:beam}.

\subsection{Candidate Generation via E-Graph Anti-Unification}

The goal of candidate generation is to find a set of patterns
that are useful for compression.
%
Following the discussion in \autoref{sec:au},
we restrict our attention to patterns $\gau(t_1, t_2)$,
where $t_1$ and $t_2$ are subterms of some $t \in \denot{\mathcal{G}}$.
%
The na\"ive approach is to enumerate all such $t$,
and for each one, perform anti-unification on all pairs of subterms;
this is suboptimal at best,
and impossible at worst (when $\denot{\mathcal{G}}$ is infinite).
%
Hence in this section we show how to compute a finite set of candidate patterns 
directly on the e-graph,
without discarding any optimal patterns.

\mypara{E-Class Anti-Unification}
%
Let us first consider anti-unification of two e-classes, $\gau(a,b)$,
which takes as input e-class ids $a$ and $b$
and returns a set of patterns.
%
We define $\gau(a,b) = \ecau{\emptyseq}{a}{b}$,
where $\ecau{\Gamma}{a}{b}$ is an auxiliary function
that additionally takes into account a \emph{context} $\Gamma$.
%
A context is a list of pairs of e-classes that have been visited while computing the AU,
and is required to prevent infinite recursion in case of cycles in the e-graph.


\begin{figure}
  \textbf{E-node Anti-Unification} $\ecau{\Gamma}{n_1}{n_2}$
  $$
  \begin{array}{rl}
    \ecau{\Gamma, (a,b)}{s(a_1,\ldots,a_k)}{s(b_1,\ldots,b_k)} &= \{s(p_1,\ldots,p_k) \mid p_i \in \ecau{\Gamma}{a_i}{b_i}\}\\
    \ecau{\Gamma, (a,b)}{s_1(\ldots)}{s_2(\ldots)} &= \{X_{a,b}\} \quad\text{if}\ s_1 \neq s_2\\
    \end{array}
  $$
  \textbf{E-class Anti-Unification} $\ecau{\Gamma}{a}{b}$
  $$
  \begin{array}{rl}
    \ecau{\Gamma}{a}{b} &= \emptyset \quad\text{if}\ (a,b) \in \Gamma\\
    % \textcolor{darkblue}{\ecau{\Gamma}{a}{a}} &\textcolor{darkblue}{= \extract_{\sz}(a)} \\
    \ecau{\Gamma}{a}{b} &= \textcolor{darkblue}{\dominant} \left(
       \bigcup_{n_a \in M(a), n_b \in M(b)} \ecau{\Gamma, (a, b)}{n_a}{n_b} 
    \right)
  \end{array}
  $$
\caption{E-class anti-unification defined as two mutually-recursive functions of e-nodes and e-class ids.}\label{fig:ecau}
\end{figure}

\autoref{fig:ecau} defines this operation using two mutually recursive functions
that anti-unify e-classes and e-nodes.
%
Note that e-node AU is always invoked in a non-empty context.
%
The first equation anti-unifies two e-nodes with the same constructor:
in this case, we recursively anti-unify their child e-classes
and return the cross-product of the results.
%
The second equation applies to e-nodes with different constructors:
as in term AU, this results in a pattern variable.
%
A nice side-effect of dealing with e-graphs
% is that we can use e-class ids (instead of full terms) 
% in the names the pattern variables $X_{a,b}$.
is that we need not keep track of the anti-substitution
to guarantee that each pair of subterms maps to the same variable:
because any duplicate terms are represented by the same e-class,
we can simply use the e-class ids $a$ and $b$ in the name of the pattern variable $X_{a,b}$.
%
% Although this causes the algorithm to generate equivalent patterns
% (when a pattern has more than one match in either of the e-classes),
% they can be be de-duplicated after the fact.

The second block of equations defines anti-unification of e-classes
(let us first ignore the \textcolor{darkblue}{\sf dominant} which will be explained shortly).
%
The first equation applies when $a$ and $b$ have already been visited:
in this case, we break the cycle and return the empty set.
%
Otherwise, the last equation anti-unifies all pairs of e-nodes from the two e-classes
in an updated context and merges the results.
%
Note that this will add a pattern variable unless all e-nodes in both e-classes
have the same constructor;
we have omitted this detail in \autoref{sec:overview} for simplicity,
but this is implemented in \tool and often yields more optimal patterns.
%
For example, consider anti-unifying $c_1$ and $c_2$ from \autoref{fig:egraph} (right):
although they have the constructor \T{scale} in common,
$X$ is actually a better pattern for abstracting these two e-classes than $\T{scale}\  X\ Y$,
because the total size of the actual arguments to pattern $X$ ($\pSide{6}$ and $\scale{\repRot{\pSide{8}}{8}{2\pi/8}}{2}$)
is the same as those to the pattern $\T{scale}\  X\ Y$ ($\pSide{6}$, $1$, $\repRot{\pSide{8}}{8}{2\pi/8}$, and $2$),
and the pattern $X$ itself is smaller.
%
This happens because the class $c_1$ represents several different terms, 
and the term ``compatible with'' $X$ in this case is smaller than the term ``compatible wit'' $\T{scale}\  X\ Y$.

%  because the argument to $X$ (\ie \T{(side 6)}) is shorter than 
%  the total size of both arguments to $Y$ (\ie \T{(side 6)} and \T{1}).

\mypara{Dominant Patterns}
% \mw{this part is a bit confusing, could come back to it}
% \mw{
%   What are we claiming here? That we aren't throwing away anything optimal?
%   Is it possible that a dominated pattern actually matches more terms, and thus is better?
% }
%
Recall that $\gau(a,b)$ produces patterns
with variable names $X_{a_i,b_i}$, which record the e-class ids they abstract;
let us refer to such a pattern $p$ as \emph{uniquely matched}
and define $\theta_l(p) = \{\many{X_{a_i,b_i} \mapsto a_i}\}$ and $\theta_r(p) = \{\many{X_{a_i,b_i} \mapsto b_i}\}$;
these substitutions are necessarily among $\matches{a}{p}$ and $\matches{b}{p}$, respectively.
%
Given two uniquely matched patterns $p_1,p_2 \in \gau(a,b)$,
we say that $p_1$ \emph{dominates} $p_2$ in the context of $(a,b)$ if
  (1) $\fv(p_1) \subseteq \fv(p_2)$, and
  (2) $\sz(p_1) \leq \sz(p_2)$.
%
We can show that if $p_1$ dominates $p_2$,
then we can safely discard $p_2$ from the set of candidate patterns.
%
First, since $p_1$ is no larger than $p_2$, 
the definition of its $\lambda$-abstraction is also no larger.
%
Second, given a term $t_a\in \denot{a}$,
compressing this term using $p_1$ vs $p_2$,
requires choosing actual arguments from $\rng{\theta_l(p_1)}$ vs $\rng{\theta_l(p_2)}$;
because the former is a subset of the latter,
the first application can always be made no larger
(symmetric argument applies for $t_b\in \denot{b}$).

Hence it is sufficient that $\gau(a,b)$ only returns the set of \emph{dominant patterns}
(\ie a pattern dominated by any pattern in the set can be discarded).
%
This is what the function $\textcolor{darkblue}{\dominant}$ does in the last equation of \autoref{fig:ecau}.
%
This pruning technique is especially helpful in the presence of equational theories.
%
Suppose our theory contains the equation $X + Y \semeq Y + X$,
and that the original term $t$ contains subterms $1 + 2$ and $3 + 1$.
%
After saturation, the e-graph will represent $1 + 2$ and $2 + 1$ in some e-class $a$
and $3 + 1$ and $1 + 3$ in another e-class $b$;
$\gau(a,b)$ will then produce both patterns $X + 1$ and $1 + X$,
but it is clearly redundant to have both, 
since they match the same e-classes with the same substitutions $\theta$.
%
Pruning of dominated patterns will eliminate one of them.

\mypara{Avoiding Cycles}
%
Interestingly enough,
the same notion of dominant patterns justifies why we need not follow cycles in the e-graph when computing $\gau(a,b)$,
or, alternatively, why a finite set of candidate patterns is sufficient
to compress any term in $\denot{\mathcal{G}}$, even if this set is infinite.
%
Removing the first equation that short-circuits cycles can only lead to solutions $p'$ of the form
$$
p' = \sapp{[X \mapsto p]}{q}
$$
where $p \in \gau(a,b)$ is another solution, and $q$ is some context, added by the cycle.
%
It is clear that any such $p'$ is dominated by $p$, and hence can be discarded.

% \mypara{Self Anti-Unitification}
% %
% \TODO{I think we should remove this optimization from here and the figure: 
% seems more low-level, and the proof is non-obvious.}
% %
% The final optimization in \autoref{fig:ecau} for anti-unifying an e-class with itself.
% %
% In this case we can show that it is sufficient to simply extract the smallest term from $a$.
% %
% The argument is as follows:
% if $a$ occurs multiple times as a sub-term in the corpus,
% we can either pick the same term from it every time or pick different terms.
% %
% If we do pick the same term,
% then it is best to pick the smallest term and perform common sub-expression elimination on that term,
% which is what the pattern $\extract_{\sz}(a)$ will let us do.
% %
% Picking different terms, also cannot be more optimal,
% because in this case less than the entire term will be reused by the pattern.

% \TODO{Would be nice to add a walk-through example \eg with the e-graph from \autoref{fig:egraph}.}

\mypara{E-Graph Anti-Unification}
%
The algorithm $\gau(a,b)$ computes a set of patterns that can be used to compress terms
represented by the e-classes $a$ and $b$.
%
Our ultimate goal, however, is to compute candidate patterns for abstracting
\emph{all subterms} of some $t \in \denot{\mathcal{G}}$.
%
The most straightforward way to achieve this is to apply $\gau(a,b)$ to all pairs of e-classes in $\mathcal{G}$.
%
We can do better, however:
some pairs of e-classes need not be considered,
because they cannot occur together in a single term $t$.
%
For an example, consider the following e-graph, with $\mathcal{I} = \mathbb{N}$ and $r = 0$:
$$
0 \mapsto \{f(1),g(2)\}\quad 1\mapsto\{g(3)\} \quad 2\mapsto \{f(3)\} \quad 3\mapsto \{a\}
$$
This e-graph can result, for example, by rewriting a term $f(g(a))$
using an equation $f(g(X)) \semeq g(f(X))$.
%
In this e-graph, the e-classes $1$ and $2$ (representing $g(a)$ and $f(a)$, respectively)
clearly cannot \emph{co-occur} in the same term:
since the e-graph only represents two terms, $f(g(a))$ and $g(f(a))$.

To formalize this intuition, we define the co-occurrence relation between two e-class ids as follows.
%
An e-class $a$ is a \emph{sibling} of $b$ if there is an e-node that has both $a$ and $b$ as children.
%
An e-class $a$ is an \emph{ancestor} of $b$ if $a = b$ or $b$ is a child of some e-node $n \in M(c)$ and $a$ is an ancestor of $c$;
$a$ is a \emph{proper ancestor} of $b$ if $a$ is an ancestor of $b$ and $a \neq b$.
%
Two e-classes $a$ and $b$ are \emph{co-occurring} if
  (1) one of them is a proper ancestor of another, or
  (2) they have ancestors that are siblings.

Finally, to compute the set $\gau(\mathcal{G})$ of all candidate patterns for an e-graph $\mathcal{G}$,
\T{LLMT} first computes the co-occurrence relation between all e-classes in $\mathcal{G}$,
and then computes $\gau(a,b)$ of all pairs of e-classes that are co-occurring.
%
As mentioned in \autoref{sec:overview},
we use a dynamic programming algorithm that memoizes the results of $\gau(a,b)$
to avoid recomputation.

% Note that this will add a pattern variable unless all e-nodes in both e-classes
% have the same constructor;
% we have omitted this detail in \autoref{sec:overview} for simplicity,
% but this is required for completeness and is implemented in \tool.
% %
% For example, consider anti-unifying $c_1$ and $c_2$ from \autoref{fig:egraph} (right):
% although they have the constructor \T{xform} in common,
% $X$ is actually a better pattern for abstracting these two e-classes than $\T{xform}\  X\ Y$,
% because the term \T{line} is shorter than \T{xform line 1}.


